\documentclass[12pt]{article}
\usepackage{seantex}

% --- Document Info ---
\title{Analysis II: Woche 10}
\author{Sean Findlay}
\date{\today}

% --- Start Document ---
\begin{document}

\maketitle

\begin{definition}{$C^k$ domain}{}
    Fiven $k \geq 1$, we say that $\Omega \subset \R^n$ open and bounded is a \textbf{$C^k$ domain}, if $\forall p \in \partial \Omega\ \exists r > 0$ and $F$ rigid motion, we have 
    \begin{itemize}
        \item $F(B_r(p)) = B_r(0)$
        \item $F(\Omega) \cap B_r(0)$ satisfies: $\Omega \subset \R^n$ open, $V \subset \R^{n-1}$ open, $g: V \rightarrow \R$, $C^1$ function,
        $$
            \partial \Omega \cap B_r(0) = \left\{ (x', x_n) \in V \times \R  \ |\ x_n = g(x')\right\} \cap B_r(0)
        $$

        $$
            \Omega \cap B_r(0) = \left\{(x', x_n) \in V \times \R  \ |\ x_n < g(x') \right\} \cap B_r(0)
        $$
    \end{itemize}
\end{definition}

\begin{theorem}{Divergence/Gauss/Green/Stokes theorem}{}
    Let $\Omega \subset \R^n$ be a bounded $C^1$ domain. Let $G \in C^1(\overline\Omega, \R^n)$ (the "vector field"). Then
    $$
        \int_\Omega \mathrm{div} G \dd x = \int_{\partial \Omega} G \cdot \nu \dd \mathrm{vol_{\partial \Omega}}
    $$
\end{theorem}

\underline{Basic idea}: Take function we want to integrate, split function using partition of unity, take integral over each piece and then sum these integrals.

The divergence is: $\mathrm{div} F := \sum_{i = 1}^{N} \partial_i F_i$.

\begin{proofbox}
    By definition of a $C^1$ domain, $\forall p \in \partial \Omega,\ \exists r > 0$ and $F$, such that the properties in the above definition hold. Since $\partial \Omega$ is compact (closed and bounded subset of $\R^n$), $\exists p_1,\dots,p_n,\ r_1,\dots,r_N$ such that $\partial\Omega \subset \bigcup_{l = 1}^N B_{r_l}(p_l)$. Call $F_l(x) := R_l x + a_l$, where $R_l$ is an orthogonal matrix, $a_l$ is a translation (so that $F_l$ defines a rigid motion).

    Let $\eta_l$ be the corresponding partition of unity associated with these balls (as in the lemma above on partition of unity), with $\tilde\eta$ denoting the background function. Define
    $$
        G_l := GF \cdot \eta_l,\ \tilde G := G \tilde\eta
    $$

    Notice that
    $$
        \sum_{l = 1}^N G_l + \tilde G = G \quad (\text{in } \Omega)
    $$

    And $\mathrm{spt}(G_l) \subset B_{r_l}(p_l)$, $\mathrm{spt}(\tilde G) \subset \Omega$.

    By construction, up to applying the isometry $F_l$, $\Omega \cap B_{r_l}(p_l)$ satisfies the definition of a $C^1$ domain. So, we can apply the local divergence theorem to $G_l$ (actually, we apply it to $G_l \circ F_l^{-1}$), and thus
    $$
        \int_{\Omega} \mathrm{div} G_l \dd x = \int_{\partial \Omega} G_l \cdot \nu \dd \mathrm{vol}_{\partial \Omega}
    $$

    Moreover, since $\tilde G \in C^1(\Omega, \R^n)$, and $\mathrm{spt}(\tilde G) \subset \Omega$, we claim that
    $$
        \int_\Omega \mathrm{div} \tilde G = 0
    $$

    and assume this holds for now. We will prove this later. Then, because the divergence is a linear operator (we can see this easily),
    $$
    \begin{aligned}
        &\sum_{l = 1}^N \left( \int_{\Omega} \mathrm{div} G_l \dd x \right) + \int_\Omega \mathrm{div} \tilde G = \int_{\Omega} \left( \sum_{l = 1}^N \mathrm{div} G_l \right) \dd x  + \int_\Omega \mathrm{div} \tilde G \\
        & = \int_{\Omega} \left( \sum_{l = 1}^N \mathrm{div} G_l \right) + \mathrm{div} \tilde G \dd x = \int_\Omega \mathrm{div}(\sum_{l = 1}^{N} G_l + \tilde G) \\
        &= \int_{\partial \Omega} \left( \sum_{l = 1}^{N} G_l \right) \cdot \nu \dd \mathrm{vol}_{\partial \Omega}
    \end{aligned}
    $$
\end{proofbox}

\begin{proofbox}[of the claim]
    We will now prove that
    $$
        \int_\Omega \mathrm{div} \tilde G = 0
    $$

    We will show this by proving
    $$
        \int_\Omega \partial_i \tilde G_i \dd x = 0 \quad \forall i = 1,\dots,n
    $$

    which is equal to the divergence. Notice
    $$
        \int_\Omega G (x + te_i) \dd x = \int_{\Omega - te_i} G(x) \dd x \quad \forall t \in {-\eps, \eps}
    $$

    Here, we can apply differentiation under the integral sign:
    $$
        \implies \frac{\dd}{\dd t} \int_\Omega \tilde G(x + te_i) \dd x = \int_\Omega \partial_i G_i(x) \dd x = 0
    $$

    Alternatively, this statement can be proven using Fubini.
\end{proofbox}

\subsection{Applications of the divergence theorem}

\textbf{Archimedes principle}: (Physics) We have a body submerged in water, it will experience an upward force. Important fact from fluid dynamics: The force on a body in a fluid is always normal to the area and is $\vec F = A \cdot p \cdot \vec{\nu}$, where $A$ is the area and $p$ is the pressure, which is $p = z \cdot \rho$, with $\rho$ the density and $z$ the depth under water. So we have
$$
    \vec F = \int_{\partial\Omega} p \cdot \nu \dd \mathrm{vol}_{\partial\Omega} = \rho \int_{\partial\Omega} zv \dd S
$$

Notice that
$$
    \mathrm{div}(z e_3) = \partial_z (z e_3) = 1
$$

And
$$
\begin{aligned}
    F \cdot \vec e_3 &= \rho \int_{\partial \Omega} z (\nu \cdot e_3) \dd S = \rho \int_{\partial \Omega} (z \cdot e_3) \nu \dd S \overset{\text{thm.}}= \rho \int_{\Omega} \mathrm{div}(z e_3) \dd x \\
    &= \rho \int_\Omega 1 \dd x = \rho \cdot \mathrm{vol}_3(\Omega)
\end{aligned}
$$

which is the forula for the weight in water of the volume of $\Omega$.

\textbf{Application 2}: $B_1 \subset \R^n$. We want to measure the $(n-1)$-dimensional volume of this sphere:
$$
    \mathrm{vol}_{n-1}(\partial B_1) = \int_{\partial B_1} 1 \dd \mathrm{vol}_{n-1} = \int_{\partial \Omega} p \cdot \nu(p) \mathrm{vol}_{n-1}(p)
$$

Note that for any point on the sphere, it is a vector from the origin to the point, which is also coincidentally the normal vector $\nu$ at $x$. So we have $x = \nu(x)$.

Let $F(x) = x$. Compute the divergence:
$$
    \mathrm{div}(F) = \sum_{i = 1}^{n} \partial_i F_i = \sum_{i = 1}^{n} 1 = n
$$

So, by the divergence theorem:
$$
    \int_{\partial \Omega} p \cdot \nu(p) \mathrm{vol}_{n-1}(p) = \int_{B_1} n \dd x = \mathrm{vol}_n(B_1) \cdot n
$$

We have proven that the $(n-1)$ dimensional volume of the sphere is always the $n$-dimensional volume times $n$.


\textbf{Continuity equation}: (Fluid dynamics, general relativity, potential theory, etc.) Local conservation of relevant quantities. For example, the density of mass in a fluid. We model fluids in the following way: at every point in the fluid, the particles have a velocity $v$, and each point has a density $\rho$. The mass of fluid inside a control volume $\Omega$ is
$$
    \frac{\dd}{\dd t} \int_{\Omega} \rho \dd x = \int_\Omega \partial_t \rho \dd x = \int_\Omega \mathrm{div}(\rho v) \dd x = \int_{\partial \Omega} \rho \cdot v \cdot \vec\nu \dd S
$$

The final equation is:
$$
    \partial_t \rho = \mathrm{div}(\rho v)
$$

\textbf{Green's theorem}: $\R^2$. Let $F \in C^1(\overline\Omega, \R^2)$ be a vector field. The boundary $\partial \Omega$ is a curve. We can orient the curve by parametrising it counter-clockwise. Let $\tau(p)$ be the tangent vector to the curve at a point $p$. Green's theorem states
$$
    \int_{\partial \Omega} \vec F \cdot \tau \dd L = \int_\Omega \mathrm{rot}(\vec{F})
$$

where $\mathrm{rot}(\vec{F}) = \partial_2 F_1 - \partial_1 F_2$.

This is actually just a reformulation of the divergence theorem: Let $F = (F_1,F_2)^T$, $G = F^\perp = (-F_2, F_1)$. Apply the divergence theorem to $G$:
$$
    \int_\Omega \mathrm{div}(G) = \int_{\partial \Omega} G \cdot \nu \dd L = \int_{\partial\Omega} F \cdot \tau \dd L
$$

And the divergence $\mathrm{div}(G) = \partial_1 G_1 + \partial_2 G_2 = -\partial_1 F_2 + \partial_2 F_1$, so
$$
    \int_\Omega \mathrm{div}(G) = \int_\Omega \mathrm{rot}(\vec{F})
$$

So we have deduced from the divergence theorem that
$$
    \int_{\partial \Omega} \vec F \cdot \tau \dd L = \int_\Omega \mathrm{rot}(\vec{F})
$$

\begin{exercise}{}{}
    Show the following formulas:
    \begin{enumerate}[label=\alph*)]
        \item
        $$
            \int_\Omega \partial_i g \dd x = \int_{\partial\Omega} g \nu_i \dd \mathrm{vol_{n-1}} \quad g, h \in C^1(\overline\Omega, \R) \ \forall i = 1,\dots,n
        $$

        \item The integration by parts formula:
        $$
            \int_\Omega \partial_i g h \dd x = - \int_\Omega g \partial_i h + \int_{\partial \Omega} g h \nu_i \dd \mathrm{vol}_{n-1}
        $$

        \item $f \in C^1(\overline\Omega, \R^n)$, $g \in C^1(\overline\Omega, \R)$
        $$
            \int_\Omega \mathrm{div} f g = - \int_\Omega f \cdot \nabla g + \int_{\partial \Omega} f \cdot \nu g \dd \mathrm{vol}_{n-1}
        $$
    \end{enumerate}
\end{exercise}

\begin{notebox}{Warning}{}
    The divergence theorem also holds true in most "reasonable" domains. Consider for example a square domain in $\R^2$ (not a $C^1$ domain, it is not smooth).

    If $\Omega := \{ (x', x_n) \in \R^{n-1} \times \R \ |\ g_n(x') < x_n < g_2(x')\}$, our domain has two cusps. But the divergence theorem will still hold.
\end{notebox}

\begin{exercise}{}{}
    Prove the divergence theorem for $\Omega = [a, b] \times [c, d] \subset \R^2$.
\end{exercise}

\section{Line integrals}

Recall: Given a vector space $V$ over the real numbers, the dual space $V^*$ is defined as
$$
    V^* := \{ a: V \rightarrow \R \ |\ a \text{ is linear} \} = \mathrm{Hom}(V, \R)
$$

Take $V = \R^n$. Then, the dual space is ${\R^n}^*$. If $v \in \R^n$, we call it a vector. If $v \in {\R^n}^*$, we call it a \textbf{covector}. For our purposes, vectors are column vectors, and covectors are row vectors. (in the general case, there is more of a distinction. But for us, this is the only difference!)

\begin{definition}{}{}
    Let $U \subset \R^n$ open. A map assigning to each point $x \in U$ a covector in ${\R^n}^*$ 
    $$
    x \mapsto (\alpha_1(x),\dots,\alpha_n(x)) = \alpha_x'
    $$ 
    
    is called a \textbf{$1$-form}. It is $C^k$ if $\alpha_i: U \rightarrow \R$ of class $C^k$ (for $k \geq 1$).
\end{definition}

\begin{example}{}{}
    \begin{enumerate}[label=\alph*)]
        \item Let $U \subset \R^n$ open, $F\in C^1(U, \R^n)$ a $C^k$ vector field (a map from $U$ to $\R^n$), then
        $
            \alpha = \alpha_F \text{ defined by } \alpha_x = F(x)^T
        $ is a 1-form. (yes, it is trivial, but it is a legitimate example of a 1-form).

        \item $f \in C^k(U, \R),\ k \geq 2,\ u \ni x \mapsto D f_x$ is a $1$-form.
    \end{enumerate}
\end{example}

\begin{definition}{}{}
    $f \in C^\infty(U)$, define $df: U \rightarrow {\R^n}^*$, $df_x := Df_x$.
\end{definition}

\begin{notebox}{0-forms}{}
    We often call functions $0$-forms, $f \in C^\infty(U)$ for instance in the example above is a $0-form$.
\end{notebox}

\begin{definition}{Line integral}{}
    Let $U \subset \R^n$ open. If $\alpha: U \rightarrow {R^n}^*,\ C^1$ is a 1-form (in $U$), $\gamma: [a, b] \rightarrow U$ a $C^1$ curve, then 
    $$
        \int_\gamma \alpha := \int_{a}^{b} \alpha_{\gamma(t)} \left( \gamma'(t) \right) \dd t = \int_{a}^{b} \underbrace{\alpha_{\gamma(t)} \cdot \gamma'(t)}_{\text{matrix product}} \dd t
    $$

    Note that the matrix product is $1 \times n$ times $n \times 1$, so the result is a scalar ($1 \times 1$ matrix).
\end{definition}

\begin{lemma}{Reparametrisation invariance of the line integral}{}
    Let $\Psi: [c, d] \rightarrow [a, b]$ be a $C^1$ diffeomorphism with $\Psi' > 0$ (orientation preserving), then
    $$
        \int_\gamma \alpha = \int_{\gamma \circ \Psi} \alpha
    $$
\end{lemma}

\begin{proofbox}
    By definition,
    $$
        \int_\gamma \alpha = \int_{a}^{b} \alpha_{\gamma(t)} (\gamma'(t)),\quad 
    $$

    We can write:
    $$
    \begin{aligned}
        \int_{\gamma \circ \Psi} \alpha &= \int_{c}^{d} \alpha_{\gamma(\Psi(s))} \underbrace{ \left((\gamma \circ \Psi)'(s) \right)}_{\gamma'(\Psi(s)) \Psi'(s)} \dd s \\
        &= \int_{c}^{d} \left(\alpha_{\gamma(\Psi(s))}\right) \gamma'(\Psi(s)) \Psi'(s) \dd s
    \end{aligned}
    $$

    We can now perform a change of variables and see that the two integrals are equal.
\end{proofbox}

\begin{remark}{}{}
    To each vector field $F$, associated $\alpha = F^T$ 1-form:
    $$
        \int_{\gamma} \alpha = \int_\gamma F^T = \int_{a}^{b} F^T(\gamma(t)) \gamma'(t) \dd t = \int_{a}^{b} F(\gamma(t) )\cdot \gamma'(t) \dd t
    $$
\end{remark}

\begin{definition}{Conservative vector fields, exact 1-forms}{}
    Let $U \subset \R^n$ be a connected open set. We say that $F: U \rightarrow \R^n$ $C^1$ vector field is \textbf{conservative} if $\exists f: U \rightarrow \R,\ C^2$, such that $\nabla f = F$.
    
    Equivalently, in the language of $1$-forms, $\alpha: U \rightarrow {\R^n}^*,\ C^1$ 1-form is \textbf{exact} if $\exists f: U \rightarrow \R$ such that $\alpha = df$.
\end{definition}

Note that the transpose of the gradient is the differential (Jacobian). So the above definition makes sense.

\begin{lemma}{}{}
    If $\alpha: U \rightarrow {\R^n}^*$ is an exact 1-form, ($\alpha = df$ for some $f: U \rightarrow \R,\ C^2$), and $\gamma: [a, b] \rightarrow U \text{ a } C^1 \text{ curve}$, then
    $$
        \int_\gamma \alpha = f(\gamma(b)) - f(\gamma(a))
    $$
\end{lemma}

\begin{proofbox}
    $$
    \begin{aligned}
        \int_\gamma \alpha &= \int_{a}^{b} \alpha_{\gamma(t)} \left( \gamma'(t) \right) \dd t = \int_{a}^{b} \underbrace{df_{\gamma(t)} \left( \gamma'(t) \right)}_{\text{chain rule!}} \dd t = \int_{a}^{b} (f \circ \gamma)'(t) \dd t \\
        &\overset{\text{FTC}}= f(\gamma(b)) - f(\gamma(a))
    \end{aligned}
    $$
\end{proofbox}

\begin{lemma}{Integrability condition}{}
    Let $U \subset \R^n$ open. Let $\alpha: U \rightarrow {\R^n}^*$ be an exact $1$-form. Then, letting $\alpha_x := (\alpha_1(x),\dots,\alpha_n(x))$,
    $$
        \partial_i \alpha_i = \partial_i \alpha_j \iff \forall x \in U\ \partial_i \alpha_i(x) = \partial_i \alpha_j(x)
    $$
\end{lemma}

\begin{proofbox}
    If $\alpha = df_x = (\partial_1f(x),\dots,\partial_nf(x))$, then
    $$
        \partial_j \alpha_i = \partial_j \partial_i f \overset{\text{Schwarz}}= \partial_i \partial_j f = \partial_i \alpha_j
    $$
\end{proofbox}

When checking if a field is conservative, or if a 1-form is exact, the above lemma is the first thing to check. If this doesn't hold, the field is definitely not conservative.

\begin{lemma}{Converse of the integrability condition}{}
    Assume that $U \subset \R^n$ is convex and let $\alpha: U \rightarrow {\R^n}^*$ satisfy the integrability condition (from the lemma above). Then, $\exists f \in C^2(U)$ such that $df = \alpha$.
\end{lemma}

\begin{proofbox}
    First, prove for a box in $\R^2$. $U = (a,b)\times(c, d) \subset \R^2$. Fix $(x_0,y_0) \in U$, and fix a "destination" $(x, y) \in U$. We will travel along a path, one part horizontal then one part vertical. Start at $(x_0,y_0)$, travel via $(x, y_0)$ to $(x_0, y_0)$.
    $$
        f(x, y) = \int_{x_0}^{x} \alpha_{(t, y_0)}\left(\begin{pmatrix} 1 \\ 0 \end{pmatrix}\right) \dd t + \int_{y_0}^{y} \alpha_{(x, s)}\left(\begin{pmatrix} 0 \\ 1 \end{pmatrix}\right) \dd s
    $$
    If $\alpha_{(x, y)} = (\alpha_1(x, y), \alpha_2(x, y))$, then
    $$
        = \int_{x_0}^{x} \alpha_1(t, y_0) \dd t + \int_{y_0}^{y} \alpha_2(x, s) \dd s
    $$

    Claim: $df = \alpha \iff \nabla f = a^T$.
    \begin{proofbox}[of claim]
        $$
            \partial_1 f = \alpha_1(x, y_0) + \underbrace{\int_{y_0}^{y} \partial_1 \alpha_2(x, s) \dd s}_{\text{diff. under integral sign}} \overset{\text{int. cond.}}= \alpha_1(x, y_0) + \int_{y_0}^{yes} \partial_2 \alpha_1(x, s) \dd s
        $$
        $$
            \partial_2 f = \alpha_2(x, y) 
        $$

        $$
            \partial_1 f(x, y) = \alpha_1(x, y_0) + \alpha_1(x, y) - \alpha_1(x, y_0) = \alpha_1(x, y)
        $$
    \end{proofbox}
\end{proofbox}

\textbf{How do we generalise line integrals?} Which quantities integrated over parametrised submanifolds are \underline{invariant} under orientation preserving reparametrisation?

Let $V \subset \R^k$, $\Phi: V \rightarrow \R^n$ be a parametrised submanifold. Let $M= \Phi(V)$ (with some orientation). Then
\[
    \int_M \alpha = \int_V \alpha \left( \Phi(x), D\Phi(x) \right) \dd x
\]

where $\alpha_p(T) = \alpha(p, T)$ is a reorientation?

Let $\Psi: W \rightarrow V$ be a $C^1$ diffeomorphism with positive orientation, orientation preserving (this means $\det(D \Psi) > 0$, note that for a diffeomorphism the determinant is never zero anyway, so WLOG if $\Psi$ were to have negative orientation, then just flip the sign of one coordinate to make it positive). Let $\Phi \circ \Psi: W \rightarrow \R^n$ be another parametrisation of $M$. Then
\[
    \int_W \alpha((\Phi \circ \Psi)(y), D(\Phi \circ \Psi)(y)) \dd y = \int_W \alpha(\Phi(\underbrace{\Psi(y)}_{x}), D\Phi(\underbrace{\Psi(y)}_{x})D\Psi(y)) \dd y
\]

If we had that $\alpha(p, TS) = \alpha(p, T) \det(S)$, then
\[
    \int_{\underbrace{W}_{\Psi^{-1}(V)}} \alpha(\Phi(\underbrace{\Psi(y)}_{x}), D\Phi(\underbrace{\Psi(y)}_{x})) \cdot \underbrace{\det(D\Psi(y)) \dd y}_{\dd x}
\]

The property $\alpha(p, TS) = \alpha(p, T) \det(S)$ is given for a \textbf{$k$-form}.

If we have a fixed $p$, $T \in M_{n \times k}(\R)$, $S \in M_{k \times k}(\R)$, $TS \in M_{n \times k}(\R)$,
\[
    \alpha(TS) = \alpha(T) \det(S)
\]

If $n = k$, then for two square matrices, the determinant satisfies this. So let $\alpha(T) = \det(T)$. But if $n \neq k$, i.e. if $T$ is not a square matrix, the determinant no longer works like this. The Gram determinant would work, but only up to a sign:
\[
    \alpha(T) = \sqrt{\det T^TT} \quad\text{ then }\quad \alpha(TS) = \alpha(T) |\det(S)|
\]

One of the simplest examples for $\alpha$ that actually works is the following. Fix $A \in M_{k \times n}(\R)$, let $\boxed{\alpha_A = \det(\underbrace{AT}_{k \times k})}$. Then,
\[
    \alpha_A(TS) = \det(ATS) = \det(AT)\det(S) = \alpha_A(T)\det(S)
\]

This function $\alpha_A$ is actually part of a basis for $k$-forms, as we will see later.

Fix $A_1,\dots,A_N \in  M_{n \times k}(\R)$, let $c_1,\dots,c_N \in \R$, let 
\[
    \alpha(T) = \sum_{l = 1}^{N} c_l \det(A_l T)
\]

This definition also works, for instance.

\begin{definition}{$k$-forms in $\R^n$}{}
    A map $p \in U \subset \R^n$ assigns $\alpha_p = \alpha_p(T)$
    \begin{align*}
        \alpha_p: \R^{n \times k} &\rightarrow \R \\
        (v_1,\dots,v_k) &\mapsto \alpha_p(v_1,\dots,v_k) 
    \end{align*}

    and satisfies:
    \begin{enumerate}
        \item $\alpha_p(TS) = \alpha_p(T)\det(S) \quad \forall T,\ \forall S \in \R^{k \times k}\quad$ (the useful property we found above)
        \item $\alpha_p$ is multilinear (i.e. linear in each $v_i$)
    \end{enumerate}

    If $\alpha$ satisfies (1), (2), then
    \begin{align*}
        \alpha(T) = \sum_{1 < i_1< \dots < i_k < n} c_{i_1,\dots,i_k} \det(e_{i_1,\dots,i_k}^T T)
    \end{align*}

    where $\mathcal{I}_k = \{ (i_1,\dots,i_k)\ |\ i_1 < \dots < i_n \} \subset \{1,\dots,n\}^k$ ($\mathcal{I}_k$ is the set of all multiindices of length $k$ with elements in $\{1,\dots,n\}$, $|\mathcal{I}_k| = \binom{n}{k} = \frac{n!}{(n-k)!k!}$).

    Also $e_{i_1,\dots,i_k}^T$ is the matrix
    \[
        \begin{bmatrix}
            e_{i_1}^T \\
            \vdots \\
            e_{i_k}^T \\
        \end{bmatrix} \in M_{k \times n}
    \]
\end{definition}

\begin{definition}{Wedge product}{}
    If $\alpha$ is a $k$-form, and $\beta$ is an $m$-form with $k + m \leq n$, if

    \[
        \alpha(T) = \det(\underbrace{A}_{k \times n}\underbrace{T}_{n \times k}),\quad \beta(S) = \det(\underbrace{B}_{m \times n}\underbrace{S}_{n \times m})
    \]

    Then,
    \[
        \alpha \land \beta (X) = \det\!\left(
            \begin{bmatrix}
            A \\
            B
            \end{bmatrix}
            X
            \right)
    \]
\end{definition}

A $k$-form is called \textbf{decomposable} if we can write it as $\alpha(T) = \det(A T)$ for some fixed matrix $A$.

Note that not every $k$-form and not every $m$-form is decomposable. We can extend the definition to every $k$-form and every $m$-form in a bilinear fashion:
\[
    \left(\sum_{i = 1}^{M} a_i \alpha_i\right) \land \left(\sum_{j = 1}^{N} b_j \beta_j\right) := \sum_{i = 1}^{M} \sum_{j = 1}^{N} a_ib_j\alpha_i \land \beta_j
\]

where $\alpha_i \land \beta_j$ is defined as in the above definition (because these $\alpha$ and $\beta$ have the required form, i.e. they are decomposable).

\begin{example}{Computing the wedge product}{}
    $x \mapsto \alpha_x = (0, x_2, x_3-2x_1)$ in $\R^3$. $x \mapsto \beta_X = (1, x_3, \cos x_2)$. Then,
    \[
        (\alpha \land \beta)_x (v_1,v_2) = \det(\begin{bmatrix}
            0 & x_2 & x_3-2x_1\\
            1 & x_3 & \cos x_2 \\
        \end{bmatrix}
        \begin{bmatrix}
            v_{11} & v_{21} \\
            v_{12} & v_{22} \\
            v_{13} & v_{23} \\
        \end{bmatrix})
    \]
\end{example}

\begin{notebox}{Convention/Definition}{}
    \begin{itemize}
        \item $0$-forms in $U \subset \R^n$ open are $C^\infty(U)$ functions. (recall that a $0$-form is just a function).
        \item $k$-forms in $U \subset \R^n$,
        \[
            \sum_{I \in \mathcal{I}_k} \underbrace{c_I(x)}_{C^\infty} \det(e_I^T T) = \alpha_x(T)
        \]

        where $\mathcal{I}_k$ is the set of all increasing multiindices of length $k$ with $n$ elements. This notation allows us to write all of these linear combinations efficiently.

        \item For $k > n$, all $k$-forms $\equiv 0$. Therefore, we always take $0 \leq k \leq n$.
        \item If $k = n$, this is the special case where $|\mathcal{I}_k| = \binom{n}{n} = 1$. So $c(x) \det(I_nT) = c(x) \det(T)$ 
        \item Let $\Omega^k(U)$ be the set of $k$-forms that are $C^\infty$ over $U$.
        \item So $\Omega^0(U) = C^\infty(U)$.
    \end{itemize}
\end{notebox}

\begin{definition}{Differential}{}
    Let $U \subset \R^n$ open. If $0 \leq k \leq n$ and $\alpha \in \Omega^k(U)$ is a $k$-form, then we can of course always write:
    \[
        \alpha_x(T) = \sum_{I \in \mathcal{I}_k} c_I(x) \underbrace{\det(e_I^T T)}_{\omega_I(T)}
    \]

    Then, the \textbf{differential} $\beta := \dd \alpha$ of $\alpha$ is a $k+1$-form, which we can write as
    \[
        \beta_x(S) = \sum_{I \in \mathcal{I}_k} \dd c_I(x) \land \omega_I
    \]
\end{definition}

If $c_I \in C^\infty(U)$ is a function, or in other words, $c_I \in \Omega^0(U)$, then, the differential of $c_i$ is just
\[
    \dd c_I = Dc_I = \left( \partial_1 c_I, \dots, \partial_n c_I, \right)
\]

After all this definition, still the only useful property of $k$-forms that we know is that we can \textbf{integrate over $k$-submanifolds}, which is \textbf{reparametrisation invariant}.

\begin{remark}{}{}
    Take the $\R^n$ coordinate functions $f(x) = x_i$ for $i = \{1,\dots,n\}$. Then, the differential (Jacobian matrix) is
    \[
        \dd x_i = \dd f_x = Df(x) = e_i^T
    \]

    So, with this notation, we have found that
    \[
        \boxed{e_i^T = \dd x_i}
    \]

    So any $k$-form is of the form
    \[
        \alpha(T) = \sum_{I \in \mathcal{I}_k} c_I(x) \overbrace{\det(e_i^T T)}^{\omega_I(T)}
    \]

    We have
    \[
        \omega_I = \dd x_{i_1} \land \dots \land \dd x_{i_k} \quad \text{with } I = (i_1,i_2,\dots,i_k),\ 1 \leq i_1 < i_2 < \dots < i_k \leq n
    \]

    We also write for this (for $I \in \mathcal{I}_k$):
    \[
        \omega_I = \dd x_{i_1} \land \dots \land \dd x_{i_k} = \boxed{\dd x_I}
    \]

    With this observation, any $k$-form in $U$ has the expansion:
    \[
        \boxed{\alpha = \sum_{I \in \mathcal{I}_k} c_I \dd x_i}
    \]

    This is standard notation.
\end{remark}

\begin{exercise}{}{}
    Let $\psi: V \rightarrow U$ be a $C^\infty$ diffeomorphism with $V \subset \R^n,\ U \subset \R^n$. Denote
    \[
        \psi = \begin{pmatrix} \psi_1(x) \\ \vdots \\ \psi_n(x) \end{pmatrix}
    \]

    Show that
    \[
        \dd \psi_1 \land \dots \land \dd \psi_n = \det D \psi(x) \dd x_1 \land \dots \land \dd x_n
    \]

    Hint: show first for $n=1$, then $n=2$.
\end{exercise}

\begin{lemma}{}{}
    For every $\alpha \in \Omega^k(U)$ for $U \subset \R^n$ open. $\dd \alpha \in \Omega^{k+1}(U)$, $\dd (\dd \alpha) \in \Omega^{k+2}(U)$. Then $\dd (\dd \alpha) = 0$. In other words,
    \[
        \dd^2 = 0
    \]
\end{lemma}

\begin{proofbox}
    Take $\dd \dd (f(x) \omega_I)$.

    For the first differential, we have
    \[
        \dd f \omega_I = \sum_{i = 1}^{n} \partial_i f \dd x_i \land \omega_I
    \]

    For the second, we have
    \[
        \dd (\dd f) = \sum_{i = 1}^{n} \sum_{j = 1}^{n} \partial_j\partial_i f \underbrace{\dd x_j \land \dd x_i}_{- \dd x_i \land \dd x_j} \land \omega_I
    \]

    Next, apply Schwarz.
\end{proofbox}

\end{document}